\subsection{\bt{Published memories on RelBench}{RelBench上的已发表记忆方法对照}}
\label{sec:relbench-eval}
\bi{\emph{Earlier gateway comparison.} Table~\ref{tab:relbench} and Figure~\ref{fig:relbench-cumulative} use the archived seven-method study: 210 sessions over the six tasks below, with five cutoffs per database ending at the official RelBench split. Cost columns cover the twelve reuse sessions at cutoffs 2--5; test scores average all five cutoffs, including the entity-reference fallback for invalid solutions. This study uses 40-call and eight-validation limits, with two consecutive non-improvements ending a session and at most five further calls after the previous-solution target is reached. It includes the three memories described below, DeepPrep's official pipeline-search code~\cite{fan2026deepprep}, and our implementation of FeatPilot's time-bounded feature search~\cite{liang2025featpilot}. Reproduction details are archived separately from the subsequent official-API study.}
{\emph{较早的网关对照。}表~\ref{tab:relbench}与图~\ref{fig:relbench-cumulative}使用归档的七方法实验：在下述六项任务上共运行210个会话，每库五个截点，末截点为RelBench官方划分。代价列覆盖截点2--5的十二个复用会话；测试分数取全部五个截点的均值，无有效方案时包含实体表参照回退。本实验上限为40次调用、8次验证；连续两次不改进即停止，达到上一方案目标后最多再调用5次。比较包含下述三种记忆方法、DeepPrep官方流水线搜索代码~\cite{fan2026deepprep}，以及我们按论文实现、带时间边界的FeatPilot特征搜索~\cite{liang2025featpilot}。复现细节与后续官方API实验分别归档。}

\bi{\emph{Historical cost accounting.} Observed time sums agent-session wall-clock durations, including API latency. The USD column is a normalized estimate from recorded usage, using the archived October 22, 2026 rates per million tokens: \$0.003 cached input, \$0.15 uncached input, and \$0.60 output. These historical estimates and five-cutoff results are not pooled with the later three-cutoff official-API measurements.}
{\emph{历史代价计量。}观测时间汇总智能体会话的墙钟耗时，包含API延迟。美元列依据记录用量归一化估算，沿用2026年10月22日归档的每百万token费率：缓存输入0.003美元、未缓存输入0.15美元、输出0.60美元。历史估价与五截点结果不并入后续三截点的官方API测量。}

\begin{table*}[t]
\caption{\bt{RelBench: predictive quality, valid solutions, and reuse cost. Each method runs three tasks at five cutoffs per database.}{RelBench：预测质量、有效方案与复用代价。每种方法在每个数据库上运行三个任务、五个截点。}}
\label{tab:relbench}
\centering
\ifbilingual\footnotesize\else\small\fi
\setlength{\tabcolsep}{2.8pt}
\renewcommand{\arraystretch}{1.08}
\newcommand{\RbPanelHeader}[4]{%
  \multicolumn{12}{@{}l}{\textbf{#1}}\\[2pt]
  & \multicolumn{6}{c}{\bh{Reuse cost: cutoffs 2--5}{复用代价：截点 2--5}}
  & \multicolumn{3}{c}{\bh{Mean test score: cutoffs 1--5}{平均测试分数：截点 1--5}}
  & \multicolumn{2}{c}{\bh{Quality summary}{质量汇总}}\\
  \cmidrule(lr){2-7}\cmidrule(lr){8-10}\cmidrule(l){11-12}
  \bh{Method}{方法} & \bh{Tokens\\(M)}{token\\（百万）} & \bh{Token\\ratio}{token\\比率}
  & \bh{LLM\\calls}{LLM\\调用} & \bh{Tool\\calls}{工具\\调用}
  & \bh{Obs.\\min}{观测\\分钟} & \bh{Cost\\(USD)}{美元}
  & #2 & #3 & #4 & \bh{Valid\\sessions}{有效会话} & \bh{Quality\\tasks}{达标任务}\\
  \midrule
}
\begin{tabular*}{\textwidth}{@{\extracolsep{\fill}}lrrrrrrrrrcc@{}}
\toprule
\RbTableBody
\bottomrule
\end{tabular*}
\par\smallskip
\parbox{\textwidth}{\footnotesize\raggedright
\bt{Cost columns total 12 sessions at cutoffs 2--5, including memory operations; token ratios are relative to no memory. Dashes: not applicable (FeatPilot uses no LLM). Test scores average all five cutoffs, using the entity-table reference for sessions without a valid solution. \emph{Valid sessions}: non-fallback solutions (/15). \emph{Quality tasks}: tasks retaining no-memory predictive quality (/3), with $\mathrm{AUC}\geq\mathrm{AUC}_0-0.01$ or $\mathrm{MAE}\leq 1.02\,\mathrm{MAE}_0$. Subscript 0 denotes no memory; best task means are bold. Time and USD accounting are specified in Section~\ref{sec:relbench-eval}.}{代价列合计截点 2--5 的 12 个会话，含记忆操作；token 比率以无记忆为参照。短横：不适用（FeatPilot 不使用 LLM）。测试分数取五个截点的平均，无有效方案的会话采用实体表参照分数。有效会话计数排除参照回退（共 15 个）；达标任务指保持无记忆预测质量的任务数（共 3 项），满足 $\mathrm{AUC}\geq\mathrm{AUC}_0-0.01$ 或 $\mathrm{MAE}\leq 1.02\,\mathrm{MAE}_0$。下标 0 表示无记忆；各任务最佳平均分加粗。时间与美元计量见第~\ref{sec:relbench-eval} 节。}
}
\end{table*}

\begin{figure}[t]
\centering
% Vector drawing generated by tools/figures/relbench_cumulative.py from exp/2026-10-20-relbench/paper-sessions.json; rerun it after the numbers change.
\ifbilingual
\includegraphics[width=\linewidth]{figures/relbench-cumulative-zh}
\else
\includegraphics[width=\linewidth]{figures/relbench-cumulative}
\fi
\caption{\bt{Cumulative tokens, including memory operations, for the six LLM-based methods. Panels share the session axis; shading marks cutoff 1 (learning), followed by cutoffs 2--5 (reuse). Thin rules separate cutoffs.}{六种使用 LLM 的方法的累计 token，含记忆操作。两个面板共用会话轴；阴影标出首截点（学习），其后为截点 2--5（复用）。细线分隔各截点。}}
\label{fig:relbench-cumulative}
\Description{Two vertically stacked line charts, rel-stack and rel-trial, show cumulative tokens in millions over 15 sessions, with one shared legend. The first three sessions are shaded as the initial learning cutoff; sessions four through fifteen are subsequent cutoffs in which memory can be reused.}
\end{figure}


\bi{\emph{Five-cutoff results.} In this study, \system\ reduces reuse-stage tokens relative to no memory by \RbStackMavraTokSave\% on rel-stack and \RbTrialMavraTokSave\% on rel-trial. Including initial learning and memory operations over all five cutoffs, the reductions are \RbStackMavraTokSaveAll\% and \RbTrialMavraTokSaveAll\%. All six tasks meet the quality margins in Table~\ref{tab:relbench}, evaluated on task means across all five cutoffs. These are sequence-level results; the five temporal cutoffs are not independent experimental repetitions.}
{\emph{五截点结果。}本实验中，\system\ 相对无记忆在 rel-stack 与 rel-trial 上分别减少 \RbStackMavraTokSave\% 与 \RbTrialMavraTokSave\% 的复用阶段 token。计入全部五个截点的初始学习与记忆操作后，分别减少 \RbStackMavraTokSaveAll\% 与 \RbTrialMavraTokSaveAll\%。按各任务在五个截点上的均值评估，六项任务均达到表~\ref{tab:relbench} 的质量标准。这些结果以时间序列为单位，五个时间截点并非五次独立实验重复。}

\bi{\emph{Historical cumulative cost.} Figure~\ref{fig:relbench-cumulative} retains the original five-cutoff results and includes memory operations from initial learning onward. It plots the six LLM-based methods; FeatPilot uses no LLM and remains in Table~\ref{tab:relbench}. Across all five cutoffs, \system\ uses \RbStackMavraTokSaveAll\% fewer tokens than no memory on rel-stack and \RbTrialMavraTokSaveAll\% fewer on rel-trial. The table's reuse-only reductions are \RbStackMavraTokSave\% and \RbTrialMavraTokSave\%. These results use the earlier gateway runs; the separate official-API study below supplies the agent-work figure.}
{\emph{历史累计代价。}图~\ref{fig:relbench-cumulative}保留原来的五截点结果，从初始学习起计入记忆操作。图中包含六种使用LLM的方法；FeatPilot不调用LLM，仍列于表~\ref{tab:relbench}。全部五个截点上，\system\ 相对无记忆在rel-stack与rel-trial分别减少\RbStackMavraTokSaveAll\%与\RbTrialMavraTokSaveAll\%的token；表中仅复用阶段的降幅为\RbStackMavraTokSave\%与\RbTrialMavraTokSave\%。这些结果来自较早的网关运行；下文独立的官方API实验用于智能体工作量图。}
